\subsection*{Section 5: Monotonicity and the Limsup Lower Bound}

\begin{theorem}[Monotonicity and Limsup Bound]
\label{thm:limsup}
Let $M(n)$ denote the supremum of chromatic numbers of unit sphere packings (finite or infinite) in $\mathbb{R}^n$. Then:
\begin{enumerate}
  \item \textbf{(Monotonicity)} $M(n+1)\ge M(n)$ for all $n\ge 1$.
  \item \textbf{(Limsup lower bound)}
    \[
      \limsup_{n\to\infty}\frac{M(n)-n}{n^{2/3}}\ge 1.
    \]
\end{enumerate}
\end{theorem}

\begin{remark}
The constant $1$ is approached but not necessarily achieved: the proof gives $(M(d_q)-d_q)/d_q^{2/3} \ge (1-1/q+1/q^2)^{1/3}$, which tends to $1$ from below as $q\to\infty$.  No upper bound of the form $M(n)-n = O(n^{2/3})$ is claimed here.
\end{remark}

\begin{remark}
The gap $d_{q+1}-d_q=3q^2+q+1\sim 3d_q^{2/3}$ is of the same order as the surplus $q^2-q+1\sim d_q^{2/3}$, so a naive monotonicity interpolation does \emph{not} extend the bound $M(n)\ge n+c\,n^{2/3}$ to all large $n$ from this sequence alone.
\end{remark}

\begin{proof}
\textbf{(1) Monotonicity.}
Define the embedding $\iota:\mathbb{R}^n\to\mathbb{R}^{n+1}$ by $\iota(p)=(p,0)$.  For any unit sphere packing $\mathcal{P}$ in $\mathbb{R}^n$ (finite or infinite), the image $\iota(\mathcal{P})=\{(p,0):p\in\mathcal{P}\}$ is a unit sphere packing in $\mathbb{R}^{n+1}$, because
\[
\|\iota(p)-\iota(q)\|_{\mathbb{R}^{n+1}}=\|p-q\|_{\mathbb{R}^n}
\]
for all $p,q\in\mathcal{P}$.  In particular, distances are preserved, so tangencies are preserved, and the tangency graph of $\iota(\mathcal{P})$ is isomorphic to that of $\mathcal{P}$.  Hence every chromatic number achievable by a packing in $\mathbb{R}^n$ is also achievable by a packing in $\mathbb{R}^{n+1}$, giving $M(n)\le M(n+1)$.

\textbf{(2) Limsup bound.}
By the Main Packing Bound (Section~4), for every prime power $q\ge 2$:
\[
M(d_q)\;\ge\; q^3+1 \;=\; d_q + (q^2-q+1),\qquad d_q=q^3-q^2+q.
\]
Since $d_q=q(q^2-q+1)$, we have $d_q^{2/3}=q^{2/3}(q^2-q+1)^{2/3}$, so
\[
\frac{M(d_q)-d_q}{d_q^{2/3}}\;\ge\;\frac{q^2-q+1}{d_q^{2/3}}=\frac{(q^2-q+1)^{1/3}}{q^{2/3}}=\left(1-\frac{1}{q}+\frac{1}{q^2}\right)^{1/3}.
\]
The function $f(q)=1-1/q+1/q^2$ satisfies $f'(q)=1/q^2-2/q^3=(q-2)/q^3\ge 0$ for $q\ge 2$, so $f$ is nondecreasing for $q\ge 2$ and $\lim_{q\to\infty}f(q)=1$.

Fix any $\varepsilon\in(0,1)$.  Since $f(q)\to 1$, there exists a prime power $q_0$ such that $f(q_0)>1-\varepsilon$.  By Euclid's theorem on the infinitude of primes, there exist infinitely many prime powers $q>q_0$, giving infinitely many indices $n_j=d_{q_j}\to\infty$ with
\[
\frac{M(n_j)-n_j}{n_j^{2/3}}\;\ge\;f(q_j)^{1/3}\;>\;(1-\varepsilon)^{1/3}.
\]
Therefore $\limsup_{n\to\infty}(M(n)-n)/n^{2/3}\ge (1-\varepsilon)^{1/3}$ for every $\varepsilon\in(0,1)$.  Letting $\varepsilon\to 0$ gives
\[
\limsup_{n\to\infty}\frac{M(n)-n}{n^{2/3}}\;\ge\;1.\qquad\square
\]
\end{proof}
